\clearpage
\VerdierPhysicalPageAnchor{0217}%
\begingroup\small\setlength{\parskip}{0pt}\setlength{\emergencystretch}{1.5em}
\setlength{\abovedisplayskip}{4pt plus 2pt minus 2pt}%
\setlength{\belowdisplayskip}{4pt plus 2pt minus 2pt}%

\VerdierSemanticAnchor{ch3d-3-2-11}
\medskip\noindent\textbf{3.2.11.}\quad
Let $f=\delta^n(S^n(X,Y))$ and $g=\delta^n({S'}^n(X,Y))$ be two elements
of $\operatorname{Ext}_{\mathcal A}^{n}(X,Y)$. The morphism
$f+g\in\operatorname{Ext}_{\mathcal A}^{n}(X,Y)$ is, by definition, the
composite morphism:

\VerdierSemanticAnchor{ch3d-3-2-11-sum-composite}
\[
X\xrightarrow{\binom{\operatorname{id}}{\operatorname{id}}}X\oplus X
\xrightarrow{\begin{pmatrix}f&0\\0&g\end{pmatrix}}
Y[n]\oplus Y[n]
\xrightarrow{(\operatorname{id},\operatorname{id})}Y[n].
\]

\VerdierSemanticAnchor{ch3d-3-2-11-construction}
Paragraphs (3.2.9) and (3.2.10) then make it possible, starting from the
elements $S^n(X,Y)$ and ${S'}^n(X,Y)$, to construct an element
${S''}^n(X,Y)\in E_{\mathcal A}^{n}(X,Y)$ such that
$\delta^n({S''}^n(X,Y))=f+g$. We leave it to the reader to make this
construction explicit.

\VerdierSemanticAnchor{ch3d-3-2-12}
\medskip\noindent\textbf{3.2.12.}\quad
We have therefore shown in this subsection that the groups
$\operatorname{Ext}_{\mathcal A}^{n}(X,Y)$ ($X$ and $Y$ objects of
$\mathcal A$) can be interpreted as quotients of the sets
$E_{\mathcal A}^{n}(X,Y)$ by the equivalence relation described in (3.2.2).
Moreover, the composition of extensions and their sum can be described by
means of simple operations on the elements of the
$E_{\mathcal A}^{n}(X,Y)$. We thus recover a description due to Yoneda [3].

\VerdierSemanticAnchor{ch3d-3-3}
\subsection{The category
\texorpdfstring{$\mathcal A^{\operatorname{Ext}}$}{A Ext}.}

\VerdierSemanticAnchor{ch3d-3-3-1}
\medskip\noindent\textbf{Lemma 3.3.1.}\quad
\textit{a) Let $X$ and $Y$ be two objects of $\mathcal A$, and let
$\xi\in\operatorname{Ext}_{\mathcal A}^{0}(X[p],Y[q])\simeq
\operatorname{Ext}_{\mathcal A}^{q-p}(X,Y)$, $q>p$. There exist $q-p$ exact
sequences of length $1$, $S_i^1$, $1\leq i\leq q-p$, such that the exact
sequence of length $n$,
$S_{q-p}^1*S_{q-p-1}^1*\cdots*S_1^1$, is defined (3.2.4) and such that:}

\VerdierSemanticAnchor{ch3d-3-3-1-1}
\begin{equation}
\delta^{q-p}\bigl(S_{q-p}^1*S_{q-p-1}^1*\cdots*S_1^1\bigr)[p]=\xi.
\tag{3.3.1.1}
\end{equation}

\VerdierSemanticAnchor{ch3d-3-3-1-b-factorizations}
\noindent\textit{b) Let $S^n$ and ${S'}^n$ be two elements of
$E_{\mathcal A}^{n}(X,Y)$ that can be written, using the notation of
(3.2.6), as}
\[
S^n=S_n^1*S_{n-1}^1*\cdots*S_1^1,
\]
\[
{S'}^n={S'}_n^1*{S'}_{n-1}^1*\cdots*{S'}_1^1,
\]
\noindent\textit{where the $S_i^1$ and ${S'}_i^1$, $1\leq i\leq n$, are
exact sequences of length $1$.}

\VerdierSemanticAnchor{ch3d-3-3-1-b-equivalence}
\noindent\textit{The following properties are equivalent:}

\noindent\textit{i) $\delta^n(S^n)=\delta^n({S'}^n)$.}

\noindent\textit{ii) There exist $n$ exact sequences of length $1$,
${S''}_i^1$, $1\leq i\leq n$, and $2n-2$ morphisms of $\mathcal A$,
$f_i$, $1\leq i\leq n-1$, and $g_i$, $1\leq i\leq n-1$, such that the
sequences}
\endgroup
